\noindent
% CỘT LỀ TRÁI (Khung lịch sử Newton và Giới hạn)
\begin{minipage}[t]{0.28\textwidth}
\vspace{0pt}
\begin{tcolorbox}[
  colback=white,
  colframe=stewartgreen,
  arc=0mm,
  boxrule=0.65pt,
  left=5.5pt,
  right=5.5pt,
  top=6pt,
  bottom=6pt,
  boxsep=0pt
]
{\color{stewartgreen}\fontsize{9.5pt}{12pt}\selectfont\textbf{\textsf{Newton và giới hạn}}}\par\vspace{5pt}
{\fontsize{8pt}{10.8pt}\selectfont
Isaac Newton sinh vào ngày lễ Giáng sinh năm 1642, năm mất của Galileo. Khi vào Đại học Cambridge năm 1661, Newton chưa biết nhiều về toán học, nhưng ông tiếp thu rất nhanh qua việc đọc Euclid, Descartes và tham dự các bài giảng của Isaac Barrow. Trường Cambridge phải đóng cửa do dịch hạch hoành hành từ năm 1665 đến năm 1666, và Newton trở về quê nhà để suy ngẫm về những gì mình đã học. Hai năm đó đạt năng suất kỳ diệu đến kinh ngạc, vì chính trong thời gian này ông đã thực hiện bốn khám phá lớn của đời mình: (1) biểu diễn hàm số dưới dạng tổng của chuỗi vô hạn, bao gồm định lý nhị thức; (2) các công trình về phép tính vi tích phân; (3) các định luật chuyển động và định luật vạn vật hấp dẫn; và (4) các thí nghiệm lăng kính về bản chất của ánh sáng và màu sắc. Vì e ngại tranh cãi và chỉ trích, ông ngần ngại công bố các phát minh của mình và mãi đến năm 1687, trước sự thúc giục của nhà thiên văn học Halley, Newton mới xuất bản tác phẩm \textit{Principia Mathematica}. Trong công trình này—chuyên luận khoa học vĩ đại nhất từng được viết—Newton đã trình bày phiên bản giải tích của mình và vận dụng nó để nghiên cứu cơ học, thủy động lực học, chuyển động sóng, cũng như giải thích chuyển động của các hành tinh và sao chổi.

\vspace{6pt}
Khởi nguồn của giải tích được tìm thấy trong các phép tính diện tích và thể tích của các học giả Hy Lạp cổ đại như Eudoxus và Archimedes. Mặc dù các khía cạnh của ý tưởng giới hạn đã tiềm ẩn trong “phương pháp vét cạn” của họ, Eudoxus và Archimedes chưa bao giờ phát biểu tường minh khái niệm giới hạn. Tương tự, các nhà toán học như Cavalieri, Fermat và Barrow, những bậc tiền bối trực tiếp của Newton trong việc phát triển giải tích, cũng không thực sự sử dụng giới hạn. Chính Isaac Newton là người đầu tiên bàn luận rõ ràng về giới hạn. Ông giải thích rằng ý tưởng chủ đạo đằng sau giới hạn là các đại lượng “tiến lại gần nhau hơn bất kỳ độ sai khác định trước nào”. Newton khẳng định rằng giới hạn là khái niệm nền tảng trong giải tích, nhưng nhiệm vụ làm sáng tỏ các ý tưởng của ông về giới hạn được dành lại cho các nhà toán học thế hệ sau như Cauchy.
\par}
\end{tcolorbox}
\end{minipage}%
\hfill
% CỘT CHÍNH (Nội dung lý thuyết, ví dụ và công thức)
\begin{minipage}[t]{0.685\textwidth}
\vspace{0pt}
\noindent
Nếu ta đặt $f(x) = x$ trong Quy tắc 7 và sử dụng Quy tắc 9, ta thu được một giới hạn đặc biệt tương tự cho căn thức. (Đối với căn bậc hai, phép chứng minh được phác thảo trong Bài tập 2.4.37.)

\vspace{6pt}
% KHUNG QUY TẮC 11
\begin{tcolorbox}[
  colback=white,
  colframe=stewartred,
  arc=0mm,
  boxrule=0.65pt,
  left=8pt,
  right=8pt,
  top=6pt,
  bottom=6pt,
  boxsep=0pt
]
\noindent
\textbf{\color{stewartred}11.} \quad $\displaystyle \lim_{x \to a} \sqrt[n]{x} = \sqrt[n]{a} \qquad \text{với } n \text{ là một số nguyên dương}$
\par\vspace{3pt}
\noindent\hspace*{2.1em}(Nếu $n$ chẵn, ta giả thiết rằng $a > 0$.)
\end{tcolorbox}

\vspace{6pt}
\noindent
\textbf{\textsf{\color{stewartcyan}VÍ DỤ 2}} \quad Tính các giới hạn sau và giải thích cho từng bước biến đổi.
\par\vspace{5pt}
\noindent
\textbf{(a)} $\displaystyle \lim_{x \to 5} (2x^2 - 3x + 4)$
\hfill
\textbf{(b)} $\displaystyle \lim_{x \to -2} \frac{x^3 + 2x^2 - 1}{5 - 3x}$

\vspace{6pt}
\noindent
\textbf{\textsf{\color{stewartcyan}LỜI GIẢI}}
\par\vspace{3pt}
\noindent
\textbf{(a)} \hspace{0.2em}
$\begin{aligned}[t]
\lim_{x \to 5} (2x^2 - 3x + 4) 
&= \lim_{x \to 5} (2x^2) - \lim_{x \to 5} (3x) + \lim_{x \to 5} 4 && \text{\color{stewartcyan}(theo các Quy tắc 2 và 1)} \\[2pt]
&= 2 \lim_{x \to 5} x^2 - 3 \lim_{x \to 5} x + \lim_{x \to 5} 4 && \text{\color{stewartcyan}(theo 3)} \\[2pt]
&= 2(5^2) - 3(5) + 4 && \text{\color{stewartcyan}(theo 10, 9 và 8)} \\[2pt]
&= 39
\end{aligned}$

\vspace{6pt}
\noindent
\textbf{(b)} Ta bắt đầu bằng việc áp dụng Quy tắc 5, nhưng việc áp dụng nó chỉ thực sự được đảm bảo ở bước cuối cùng khi ta thấy rằng giới hạn của cả tử số và mẫu số đều tồn tại và giới hạn của mẫu số khác 0.
\par\vspace{5pt}
\noindent
$\begin{aligned}[b]
\lim_{x \to -2} \frac{x^3 + 2x^2 - 1}{5 - 3x} 
&= \frac{\lim\limits_{x \to -2} (x^3 + 2x^2 - 1)}{\lim\limits_{x \to -2} (5 - 3x)} && \text{\color{stewartcyan}(theo Quy tắc 5)} \\[7pt]
&= \frac{\lim\limits_{x \to -2} x^3 + 2\lim\limits_{x \to -2} x^2 - \lim\limits_{x \to -2} 1}{\lim\limits_{x \to -2} 5 - 3\lim\limits_{x \to -2} x} && \text{\color{stewartcyan}(theo 1, 2 và 3)} \\[7pt]
&= \frac{(-2)^3 + 2(-2)^2 - 1}{5 - 3(-2)} && \text{\color{stewartcyan}(theo 10, 9 và 8)} \\[7pt]
&= -\frac{1}{11}
\end{aligned}$%
\hfill {\color{stewartcyan}\rule{1.15ex}{1.15ex}}

\vspace{10pt}
\noindent
{\color{stewartcyan}\rule{1.15ex}{1.15ex}}\hspace{0.5em}\textbf{\textsf{Tính giới hạn bằng phép thay trực tiếp}}
\par\vspace{4pt}
Trong Ví dụ 2(a), chúng ta đã xác định được rằng $\lim\limits_{x \to 5} f(x) = 39$, với $f(x) = 2x^2 - 3x + 4$. Hãy chú ý rằng $f(5) = 39$; nói cách khác, chúng ta có thể nhận được kết quả chính xác chỉ đơn giản bằng cách thay $5$ vào $x$. Tương tự, phép thay trực tiếp cũng mang lại kết quả đúng ở phần (b). Các hàm số trong Ví dụ 2 lần lượt là một hàm đa thức và một hàm phân thức hữu tỉ, và việc áp dụng tương tự các Quy tắc giới hạn chứng minh rằng phép thay trực tiếp luôn có hiệu lực đối với các hàm số như vậy (xem Bài tập 59 và 60). Chúng ta phát biểu tính chất này như sau.

\vspace{6pt}
% KHUNG ĐỊNH NGHĨA: DIRECT SUBSTITUTION PROPERTY
\begin{tcolorbox}[
  colback=white,
  colframe=stewartred,
  arc=0mm,
  boxrule=0.65pt,
  left=8pt,
  right=8pt,
  top=6pt,
  bottom=6pt,
  boxsep=0pt
]
\noindent
\textbf{\textsf{\color{stewartred}Tính chất thay trực tiếp}} \quad Nếu $f$ là một hàm đa thức hoặc một hàm phân thức hữu tỉ và $a$ thuộc tập xác định của $f$, thì
\vspace{-2pt}
\[
\lim_{x \to a} f(x) = f(a)
\]
\vspace{-6pt}
\end{tcolorbox}
\end{minipage}
